\chapter{Who Proved DNA Carries Hereditary Information?}

\begin{kaichang}
If you or a relative has taken a genetic test, you have probably seen a small tube like this: spit into it, mail it away, and receive a report weeks later describing ancestry, lactose tolerance, and alcohol metabolism. Paternity testing is similar: swab the inside of the cheek, send it to a laboratory, and determine whether two people are biological father and child.

Pause over the enormous leap involved: why should this work? Why can some molecule hidden in a mouthful of saliva ``represent you''? Why should it keep the records of your blood type, tongue rolling, or tendency to get drunk?

Today the answer seems obvious: DNA, of course. Everyone knows that. In the first half of the twentieth century, however, this answer was not merely unclear; almost nobody wanted to believe it. Most leading biochemists backed a different molecule: protein. This chapter follows science, experiment by experiment, from ``probably protein'' back to ``definitely DNA.'' No single experiment settled everything. An interlocking chain of evidence secured the conclusion.
\end{kaichang}

\section{A Bet Almost Everyone Got Wrong}

Return to a 1940s laboratory and describe only what could then be measured.

Biologists already knew of threadlike nuclear structures called chromosomes: before cell division, each is copied, then shared precisely between two daughter cells (Chapter 64's transport process). Mendel's pea experiments in Chapter 63 showed that hereditary information passes in discrete ``factors.'' Put the observations together, and the natural conclusion is that chromosomes hold hereditary information.

What are chromosomes made of? Chemical analysis found two main components. One was protein, familiar from Chapter 48: chains of 20 amino-acid types that fold into countless shapes. Known enzymes, antibodies, and many hormones were proteins. The other was deoxyribonucleic acid, or DNA: chains of nucleotides differing only in four nitrogenous bases---adenine (A), guanine (G), cytosine (C), and thymine (T). Only four.

Place your bet from their historical perspective: which carries hereditary information? Most scientists chose protein, with seemingly strong reasons. Information implies enormous variety; proteins have 20 ``letters'' and abundant combinations. DNA has only four, and the popular ``tetranucleotide hypothesis'' held that they simply repeat in equal proportions. Boring and monotonous, DNA could at most serve as chromosome scaffolding. This hypothesis arose from chemist Levene's early rough measurements: base proportions happened to look nearly one-to-one, and the mistaken picture of monotonous repetition entered textbooks and stayed for decades. Twenty letters versus four---was there even a contest?

We will uncover the reasoning's flaw midway through the chapter. For now, remember: in science's history, ``everyone thinks so'' has never meant ``evidence supports it.'' No experiment had directly compared the two molecules' hereditary abilities. People had merely voted for ``whichever looks more complicated.''

\section{Dead Bacteria Change Living Bacteria's Type}

The first crack appeared in London in 1928. World War I had ended not long before, and pneumonia remained a common deadly disease. Britain's Ministry of Health wanted to understand differing bacterial virulence in preparation for vaccines. Frederick Griffith was studying pneumococcus, a pneumonia-causing bacterium, as part of this work. Two types could be distinguished by eye: S (smooth) colonies were smooth and shiny, their bacteria enclosed in polysaccharide capsules that evaded mouse immune defenses and caused disease. R (rough) bacteria lacked capsules, formed rough colonies, and did not cause disease.

Griffith performed four groups of injection experiments:

\begin{table}[htbp]
\centering
\begin{tabular}{lll}
\toprule
\shortstack[l]{Material injected\\into mice} & Outcome & \shortstack[l]{Live bacteria isolated\\after death} \\
\midrule
Live R & Survived & None \\
Live S & Died & S \\
Heat-killed S & Survived & None \\
\shortstack[l]{Live R +\\heat-killed S} & Died & Live S \\
\bottomrule
\end{tabular}
\end{table}

The first three results were unsurprising. The fourth caused the explosion: live R alone left mice alive; dead S alone also left them alive; injected together, they killed mice, and living S bacteria were isolated from the bodies. Those S bacteria continued reproducing, with all descendants remaining S.

Clarify a tempting misunderstanding: dead bacteria did not ``come back to life.'' Mice injected with heat-killed S alone remained healthy. Instead, a few living R bacteria absorbed something released by dead S bacteria and acquired the ``recipe'' for a capsule. They and their descendants became S. What revived was not a bacterial individual, but information.

Griffith called the phenomenon ``transformation'' and the invisible agent the ``transforming principle.'' Now practice a question we will repeatedly ask: \textbf{what can this experiment prove}? It proves that bacteria contain a transferable, heritable substance that can permanently alter an organism's traits at the molecular level. Experiment forced ``hereditary material'' onto the stage for the first time. But it \textbf{cannot prove} which molecule the transforming principle is. Griffith himself leaned toward a protein or capsule substance. He opened the door a crack without walking through it.

\section{Eliminating the Molecular Candidates}

Oswald Avery and his Rockefeller Institute colleagues MacLeod and McCarty took up the baton. They spent more than a decade first reproducing transformation in test tubes, without mice, then applying an elegant logical method: \textbf{elimination}.

They broke open S bacteria, purified their contents, and removed candidate molecules one by one to see whether transforming ability remained:

\begin{itemize}[itemsep=2pt]
\item Protease degraded proteins---transformation still occurred.
\item RNase degraded RNA---transformation still occurred.
\item An enzyme degraded capsule polysaccharides---transformation still occurred.
\item DNase specifically cut DNA---transformation disappeared entirely.
\end{itemize}

In 1944, they published their conclusion: the transforming principle was DNA, not merely a structural component but the functioning hereditary material.

What can this prove? Under highly purified conditions, DNA itself suffices to transfer a hereditary trait from one bacterium to another. Information still functions without protein, RNA, or polysaccharides present. Was the logic airtight? Almost. Critics seized on one point: however pure an extract, who could guarantee it contained no residual protein at a few parts per million? Perhaps that trace protein did the real work. This may look like hair-splitting today, but when ``everyone backed protein,'' it kept most observers on the fence. Avery himself was exceptionally cautious, his paper restrained. Work later regarded as Nobel-worthy received only polite silence at the time. Evidence's weight, you see, depends not only on an experiment but also on whether its era is ready.

\section{Different Glowing Labels for Two Molecules}

Alfred Hershey and Martha Chase pushed the door open in 1952. They chose a cleaner system: T2 bacteriophage, a virus that infects bacteria. ``Cleaner'' meant only two components remained---a protein coat around a DNA molecule---reducing thousands of bacterial candidate molecules to two. Infection resembles an injection: the coat attaches to a bacterium and injects its ``contents.'' The bacterial factory is forced to mass-produce phages, then bursts to release hundreds or thousands of offspring.

The key question: is DNA or protein injected? Whichever enters must be the hereditary information, because the offspring's coats and DNA are manufactured according to the ``entrant's'' instructions.

Hershey and Chase's trick was to attach different radioactive labels to the molecules. Chapter 4 discussed isotopes: isotopes of one element behave identically chemically, but radioactive isotopes emit detectable radiation. The molecules conveniently differ in elemental composition. Proteins contain sulfur but almost no phosphorus (two amino acids contain sulfur); DNA's sugar--phosphate backbone contains abundant phosphorus and no sulfur. Thus:

\begin{itemize}[itemsep=2pt]
\item Grow phages in an environment containing radioactive sulfur $^{35}\mathrm{S}$---only their protein coats become labeled.
\item Grow phages with radioactive phosphorus $^{32}\mathrm{P}$---only DNA becomes labeled.
\end{itemize}

The labeled phages then infected bacteria. After several minutes, an ordinary kitchen blender vigorously shook off the empty coats attached to bacterial surfaces. Centrifugation sent bacteria to the tube bottom while the light coats stayed in the supernatant. Measuring radioactivity in each layer revealed what entered the cells:

\begin{table}[htbp]
\centering
\begin{tabular}{llll}
\toprule
Label & \shortstack[l]{Labeled\\molecule} & \shortstack[l]{Main location\\of radioactivity} & Offspring phages \\
\midrule
$^{35}\mathrm{S}$ & Protein coat & \shortstack[l]{Supernatant\\(coats)} & \shortstack[l]{No detectable\\radioactivity} \\
$^{32}\mathrm{P}$ & DNA & Bacterial pellet & \shortstack[l]{Some carry\\$^{32}\mathrm{P}$} \\
\bottomrule
\end{tabular}
\end{table}

The conclusion is clear: DNA entered the bacteria; the protein coats remained outside as ``empty syringes.'' The entering DNA directed production of hundreds or thousands of complete offspring phages, including entirely new protein coats. DNA gave the orders.

Practice again: can this experiment \textbf{alone} prove everything? No. A little $^{35}\mathrm{S}$ still sedimented with bacteria, allowing critics to suggest trace protein had slipped inside. Moreover, only one kind of organism, a phage, was studied. Its irreplaceable contribution was directly tracing each molecular type in a system with distinct ``inside'' and ``outside.'' Avery's vulnerability was possible protein left in an extract; Hershey and Chase turned that into a direct sighting of ``who enters.'' Each experiment filled gaps in the other. Textbooks pair them because evidence rests on interlocking links, not one heroic experiment.

\textbf{Translator safety note:} These pathogen, animal-injection, and radioactive-tracer experiments are historical accounts, not home or classroom instructions. Do not culture pathogens, inject animals, handle radioactive materials, or use a food blender for biological samples. Such research needs appropriate facilities, trained supervision, and formal safety review; see \href{https://institute.acs.org/acs-center/lab-safety/education-training/safer-experiments.html}{ACS safer-experiment guidance}.

\begin{wuqu}
``Hershey and Chase's blender experiment conclusively proved by itself that DNA is hereditary material.'' This is the most widespread simplification. Alone, it neither excludes trace protein entry nor extends beyond phages. The compelling case comes from complementing Avery's enzyme experiments: one shows ``purified DNA suffices to transmit traits,'' the other ``DNA enters when traits are transmitted.'' Most scientific ``single decisive blow'' stories are retrospective rewrites. The real process uses several experiments approaching one conclusion from different directions until alternative explanations are squeezed out individually.
\end{wuqu}

\section{Balancing the Four-Letter Ledger}

The phage experiment answered ``who enters,'' but an old account remained unsettled: is DNA really so monotonous? If its four bases always repeat in equal proportions, it indeed cannot store much information. Until the tetranucleotide hypothesis fell, protein's supporters had an escape route.

Biochemist Erwin Chargaff dismantled it. In the late 1940s, he used paper chromatography to separate DNA components from different organisms and measure the four bases precisely. Two findings emerged.

First, base proportions differ between species. Human, cow, bacterial, and yeast DNA all differ. DNA is not a uniform, monotonous scaffold: its ``recipe'' changes with species, a prerequisite for carrying information.

Second, stranger still: in every species, A approximately equals T and G approximately equals C. These are ``Chargaff's rules.'' His human measurements, for example, gave about 30.9\% A, 29.4\% T, 19.9\% G, and 19.8\% C. Not exact---measurements have errors---but the matching A--T and G--C amounts were no coincidence. Chargaff did not decipher the equality's meaning himself, yet it virtually shouted that DNA's bases come in \textbf{pairs}. When Watson and Crick visited in 1952, Chargaff explained these figures personally. A few years later, this equality became a key piece in assembling the double helix.

\begin{qiaoliang}
Return to the opening bet: are four letters really too few? A sequence with $n$ positions, each permitting A, T, G, or C, has $4^n$ possible forms. For $n=10$, that is about a million; for $n=100$, about $10^{60}$. A 3000-base sequence has $4^{3000}\approx 10^{1800}$ possibilities, while the observable universe contains only about $10^{80}$ atoms. Four letters offer absurdly large information capacity. The protein camp's intuition failed not through mathematics but through the tetranucleotide hypothesis: they thought DNA's letters could only repeat equally, without freedom of arrangement. Chargaff's data showed precisely that proportions vary by species and arrangement is free. Four letters are entirely enough to write a book of heredity.
\end{qiaoliang}

\section{Photo 51: Credit, Competition, and Ethics}

\begin{zhengju}
The final puzzle piece was DNA's shape. In the early 1950s, several laboratories raced to solve its three-dimensional structure. Rosalind Franklin at King's College London was among the finest X-ray diffraction experts. The principle is this: X-rays pass through an orderly molecular sample, and regularly arranged atoms scatter them into a pattern. Spot positions reveal molecular geometry, rather like inferring a fence's spacing from its patterned shadow.

In 1952, Franklin and graduate student Gosling produced the famous ``Photo 51,'' a DNA-fiber diffraction image with a clear central X. An expert immediately recognizes that X as a helix's diffraction signature. The pattern also encodes the helix's diameter and rise per turn. Franklin's measurements further showed that DNA's hydrophilic backbone lies outside the molecule.

In January 1953, her colleague Wilkins showed Photo 51 to the visiting Watson without her consent. Watson later recalled that his mouth fell open: the helical hypothesis was confirmed. Meanwhile, measurements Franklin had assembled in an internal report also reached Crick. Within weeks, Watson and Crick assembled the double-helix model: two strands winding in opposite directions, A paired with T and G with C, explaining Chargaff's equalities. Their paper appeared in April 1953, alongside separate data papers from Franklin and Wilkins.

The next part must be told carefully. Franklin had left King's College to work on tobacco mosaic virus, again producing outstanding results. She died of ovarian cancer in 1958, aged only 37. The 1962 Nobel Prize in Physiology or Medicine went to Watson, Crick, and Wilkins. Nobel Prizes go only to the living, leaving her no opportunity even for nomination. Watson's memoir The Double Helix portrayed her as a difficult supporting character who failed to understand her own data's value. Many readers learned only years later, from collected historical records, that she had produced Photo 51 and measured the crucial data.

How should we judge this history? History does not answer ``what if'': we cannot know whether another year would have allowed Franklin to solve the structure herself. But the ethical baseline is clear: \textbf{using data requires consent, and contributions should be recorded honestly}. Cooperation and competition both advance science, but ``who first saw a photo'' must not eclipse ``who produced it.'' Today almost every textbook includes Franklin alongside the double helix. This is not charity; it is a scientific community belatedly balancing its ledger.
\end{zhengju}

\section{The Evidence Chain in One Table}

Step back and place all five links together: what did each establish, and what did it leave unanswered?

\begin{table}[htbp]
\centering
\small
\begin{tabular}{p{3.0cm}p{4.8cm}p{4.8cm}}
\toprule
Experiment or discovery & What it establishes & What it cannot establish alone \\
\midrule
Griffith (1928) & A transferable, heritable ``transforming principle'' exists & Which molecule is the transforming principle \\
Avery (1944) & Highly purified DNA suffices to transmit traits & No trace impurities remain; the conclusion holds for other organisms \\
Hershey--Chase (1952) & DNA enters bacteria during phage infection & No protein enters at all; every organism uses DNA \\
Chargaff (from 1949) & Base proportions vary by species, with A$\approx$T and G$\approx$C & DNA's specific three-dimensional structure \\
Franklin and the double helix (1953) & DNA's helical geometry and base-pairing arrangement & How structure guides replication (next chapter) \\
\bottomrule
\end{tabular}
\end{table}

This table is the chapter's boundary level: individual experiments leave gaps; the whole chain closes them. It also lays the first foundation for the book's fifth recurring question: ``How is information stored, copied, regulated, and selected in life?'' Before asking how it is stored or copied, we must identify its molecular carrier. That hard-won identification took a quarter century. Draw one more boundary: ``DNA is hereditary material'' more precisely means ``DNA is the \textbf{principal} hereditary material.'' In 1956, scientists disassembled and reconstituted tobacco mosaic virus, showing that its hereditary material is RNA and needs no DNA. Influenza virus and the coronavirus responsible for COVID-19 are also RNA viruses. You may also have heard that ``mad-cow prions transmit heredity through protein.'' Prions are indeed ``infectious'' proteins, but they carry no sequence information; they induce normal proteins to adopt the same faulty fold, rather than recording unlimited varied information like DNA or RNA. Calling them ``hereditary material'' extends beyond the evidence.

\begin{shiyan}
DNA sounds sophisticated, yet you can ``pull'' it from fruit at home. Materials: half a ripe banana (or two strawberries), a small spoonful of salt, dishwashing liquid, warm water, gauze or a coffee filter, transparent cups, and chilled medical alcohol (or strong baijiu).

Steps: (1) Mash the banana in a food-storage bag, add salt water (one small spoonful in half a cup of warm water) and a few drops of dishwashing liquid, mix gently, and leave for 10 minutes. Detergent dissolves the phospholipid bilayers of the cell and nuclear membranes (Chapter 51); salt helps DNA aggregate. (2) Filter through gauze and retain the liquid. (3) Slowly pour cold alcohol down the cup's wall so it floats above the filtrate; do not stir. (4) Wait several minutes and observe the interface.

Observations: White, cloudlike strands slowly appear at the interface and can be wound around a toothpick: clumped DNA. You are seeing tangled aggregates of trillions of DNA molecules. One double helix is only 2 nanometers wide and cannot be seen with any light microscope.

Safety: Alcohol is flammable; keep it away from flames throughout. Do not eat anything that has touched detergent. Keep detergent and alcohol out of your eyes.
\end{shiyan}

\textbf{Translator safety note:} An adult should handle and dispense the alcohol, away from flames, heat, and sparks; never heat it or taste any extract. Keep alcohol and detergent away from children and pets, and use eye protection and clearly labeled containers. Medical alcohol is not a drink; see \href{https://www.poison.org/articles/rubbing-alcohol-only-looks-like-water}{Poison Control alcohol guidance} and \href{https://www.acs.org/education/celebrating-chemistry-editions/2024-ccew/science-safety-tips.html}{ACS activity precautions}. The visible precipitate is a crude mixture, not proof of purified DNA.

\section{Back to the Saliva Tube}

Now answer the opening question. Your saliva sample contains cells shed from the cheek lining. Each cell's nucleus holds 46 chromosomes wound with DNA. Half your DNA came from your father and half from your mother---Chapter 64 explained meiosis's ``equal split.'' Among billions of base pairs, about one letter in a thousand differs from other people's. These differences are everything that paternity tests, forensic comparisons, and genetic tests read. Essentially, a laboratory reads the letters and compares who resembles whom.

Feel the scale of this ``book.'' Straightened DNA from an ordinary human cell is about two meters long and contains roughly six billion base pairs. Your body has about $3\times 10^{13}$ cells, so all its DNA joined end to end would stretch $6\times 10^{13}$ meters, or sixty billion kilometers. Earth is about $1.5\times 10^{11}$ meters from the Sun; division gives 400 Earth--Sun distances, enough for 200 round trips. All this material was identified as an ``information carrier'' link by experimental link, beginning with Griffith's mouse.

That chain of evidence has long since left textbooks. DNA comparison has helped solve countless crimes and clear the innocent. Forensic laboratories do not read whole genomes, but compare a dozen to twenty variable ``short tandem repeat'' sites; a complete match between two people becomes almost impossibly unlikely. During the pandemic, nucleic acid testing became familiar to hundreds of millions. It detects a virus's own nucleic-acid sequence, confirming ``nucleic acids carry information'' by another route. Technologies for reading and rewriting DNA are becoming everyday medical and agricultural tools. Chapter 82 will discuss what ``rewriting genes'' can achieve and where it should stop.

\section{The Carrier Is Found; Its Reading Remains}

We end with ``what is hereditary material?'' settled, but a deeper question emerging: why does DNA's structure allow both \textbf{stable storage} of billions of letters and \textbf{precise copying} before cell division? Even on paper, printing and handwriting have vastly different error rates. How can molecular ``copying'' reproduce billions of letters almost flawlessly? Chargaff's A-equals-T rule and Photo 51's X have half-revealed the answer. Why do A and T pair precisely? Why do the strands wind in opposite directions? How does base pairing automatically ensure copying accuracy? Next chapter, ``Why DNA's Structure Suits Information Storage,'' we take the double helix apart and see how it serves as both ``safe'' and ``photocopier.''

\begin{tiaozhan}
Why can Hershey--Chase labels mark ``only protein'' and ``only DNA''? Element distribution is key: among the 20 amino acids, only methionine and cysteine contain sulfur, whereas DNA contains none. Every nucleotide in DNA's sugar--phosphate backbone carries phosphorus, while proteins contain almost none. Growth with $^{35}\mathrm{S}$ therefore labels only protein coats; growth with $^{32}\mathrm{P}$ labels only DNA. For detection, $^{32}\mathrm{P}$ has a half-life of about 14.3 days and $^{35}\mathrm{S}$ about 87 days. Their radiation signals a counter or photographic film, so labels remain ``lit'' throughout experiments completed within weeks. Another historical entry deserves mention: Hershey shared the 1969 Nobel Prize for this work, while Chase, who performed most of the hands-on experiments, was credited then only as a technician. Scientific credit has more than Franklin's example.
\end{tiaozhan}

\section*{Try It Yourself}

\begin{sikaoti}
\item Draw an information-flow diagram for Griffith's fourth group: where does ``capsule-making information'' begin, what carries it, and where is it ultimately expressed? Use boxes and arrows, noting that information did not disappear with dead bacteria.
\item As a 1944 critic convinced ``residual protein is the real culprit'' in Avery's experiment, what new experiment would you design to test that objection? (Hint: what did Hershey and Chase do eight years later?)
\item Why could Hershey and Chase not distinguish protein and DNA using carbon or nitrogen labels? (Hint: which elements occur in both molecular types?)
\item Chargaff measures 24\% A in an organism's DNA. Approximately what percentage is T? What is the combined percentage of G and C? (Remember that the four bases total 100\%.)
\item How many possible sequences can a 20-base DNA segment have? Approximately how many times the world's population (about $8\times 10^{9}$) is that?
\item Regarding Photo 51: what should you have done in Wilkins's position? Discuss whether ``competition'' conflicts with ``following rules.'' Then research Franklin's life and how she is commemorated today.
\end{sikaoti}
